% =====================================================================
%  minggu-04.tex  --  Slide Pertemuan Minggu ke-4
%  Pengantar Kriptografi (MPFA014) -- Prodi Matematika FMIPA UBB
%
%  Topik silabus : Kriptografi klasik I
%  Sumber materi : buku_ajar/chapters/04/section-01.tex dan section-05.tex
%  Sub-CPMK      : 1.1
%
%  Cakupan: Caesar, substitusi monoalfabetik, Affine, transposisi Rail
%           Fence, dan super-enkripsi. Vigenere, Playfair, Hill, Enigma,
%           dan OTP tidak dibahas pada pertemuan ini.
%
%  Kompilasi dari dalam folder ppt/:
%      compile-ppt.bat 04
% =====================================================================

\documentclass[aspectratio=169,11pt]{beamer}
% =====================================================================
%  beamer-preamble.tex  --  prasarana dan gaya slide kuliah
%
%  Mata kuliah : Pengantar Kriptografi (MPFA014)
%  Program     : Program Studi Matematika, FMIPA, Universitas Bale Bandung
%
%  Berkas ini TIDAK memuat \documentclass. Setiap berkas minggu memuatnya
%  sendiri lalu menarik berkas ini:
%
%      \documentclass[aspectratio=169,11pt]{beamer}
%      % =====================================================================
%  beamer-preamble.tex  --  prasarana dan gaya slide kuliah
%
%  Mata kuliah : Pengantar Kriptografi (MPFA014)
%  Program     : Program Studi Matematika, FMIPA, Universitas Bale Bandung
%
%  Berkas ini TIDAK memuat \documentclass. Setiap berkas minggu memuatnya
%  sendiri lalu menarik berkas ini:
%
%      \documentclass[aspectratio=169,11pt]{beamer}
%      % =====================================================================
%  beamer-preamble.tex  --  prasarana dan gaya slide kuliah
%
%  Mata kuliah : Pengantar Kriptografi (MPFA014)
%  Program     : Program Studi Matematika, FMIPA, Universitas Bale Bandung
%
%  Berkas ini TIDAK memuat \documentclass. Setiap berkas minggu memuatnya
%  sendiri lalu menarik berkas ini:
%
%      \documentclass[aspectratio=169,11pt]{beamer}
%      \input{beamer-preamble}
%      \title{...} \subtitle{...} \setWeekInfo{Minggu n}{Topik}
%
%  Berkas minggu hanya berisi \setWeekInfo dan lingkungan frame.
%  Kompilasi dijalankan DARI DALAM folder ppt/ (lihat compile-ppt.bat),
%  karena gambar buku dirujuk dengan jalur relatif ../buku_ajar/.
% =====================================================================

% ---------------------------------------------------------------- bahasa
\usepackage[utf8]{inputenc}
\usepackage[T1]{fontenc}
\usepackage[indonesian]{babel}
\usepackage{lmodern}
\usepackage{microtype}

% ------------------------------------------------------------- matematika
\usepackage{amsmath,amssymb,mathtools}

% ----------------------------------------------------------------- tabel
\usepackage{booktabs}
\usepackage{array}
\usepackage{tabularx}

% --------------------------------------------------------------- grafik
\usepackage{tikz}
\usetikzlibrary{arrows.meta,positioning,shapes.geometric,calc,fit,backgrounds}

% ---------------------------------------------------------- kode program
\usepackage{listings}
\lstdefinestyle{slidecode}{
    basicstyle=\ttfamily\scriptsize,
    keywordstyle=\color{ubbblue}\bfseries,
    commentstyle=\color{ubbgreen}\itshape,
    stringstyle=\color{ubbred},
    numbers=left,
    numberstyle=\tiny\color{gray},
    numbersep=6pt,
    frame=single,
    rulecolor=\color{ubbblue!40},
    backgroundcolor=\color{ubbblue!3},
    breaklines=true,
    showstringspaces=false,
    tabsize=4,
    columns=fullflexible,
    keepspaces=true,
}
\lstset{style=slidecode}

% =====================================================================
%  WARNA INSTITUSI
% =====================================================================
\definecolor{ubbblue}{RGB}{12, 62, 120}      % warna utama
\definecolor{ubbteal}{RGB}{0, 118, 110}      % aksen kedua
\definecolor{ubbgreen}{RGB}{22, 122, 74}     % untuk blok contoh
\definecolor{ubbred}{RGB}{168, 42, 42}       % untuk blok peringatan
\definecolor{ubbgray}{RGB}{70, 76, 84}       % teks sekunder

% =====================================================================
%  IDENTITAS MATA KULIAH
% =====================================================================
\newcommand{\CourseName}{Pengantar Kriptografi}
\newcommand{\CourseCode}{MPFA014}
\newcommand{\CourseSemester}{Semester Ganjil}
\newcommand{\CourseProgram}{Program Studi Matematika}
\newcommand{\CourseFaculty}{Fakultas Matematika dan Ilmu Pengetahuan Alam}
\newcommand{\CourseInstitution}{Universitas Bale Bandung}
\newcommand{\CourseLecturer}{Cecep Suwanda, S.Si., M.Kom.}

% Label minggu, disetel tiap berkas lewat \setWeekInfo
\newcommand{\WeekLabel}{}
\newcommand{\WeekTopic}{}
\newcommand{\setWeekInfo}[2]{%
    \renewcommand{\WeekLabel}{#1}%
    \renewcommand{\WeekTopic}{#2}%
}

% =====================================================================
%  TEMA: header, footer, warna blok
% =====================================================================
\usetheme{default}
\setbeamertemplate{navigation symbols}{}
\setbeamertemplate{blocks}[rounded][shadow=false]

\setbeamercolor{palette primary}{bg=ubbblue, fg=white}
\setbeamercolor{palette secondary}{bg=ubbgray, fg=white}
\setbeamercolor{palette tertiary}{bg=ubbteal, fg=white}

\setbeamercolor{structure}{fg=ubbblue}
\setbeamercolor{normal text}{fg=black!92}
\setbeamercolor{alerted text}{fg=ubbred}

\setbeamercolor{block title}{bg=ubbblue, fg=white}
\setbeamercolor{block body}{bg=ubbblue!6, fg=black!92}
\setbeamercolor{block title alerted}{bg=ubbred, fg=white}
\setbeamercolor{block body alerted}{bg=ubbred!6, fg=black!92}
\setbeamercolor{block title example}{bg=ubbgreen, fg=white}
\setbeamercolor{block body example}{bg=ubbgreen!6, fg=black!92}

\setbeamercolor{itemize item}{fg=ubbblue}
\setbeamercolor{itemize subitem}{fg=ubbteal}
\setbeamercolor{enumerate item}{fg=ubbblue}

\setbeamerfont{headline}{size=\scriptsize}
\setbeamerfont{footline}{size=\scriptsize}
\setbeamerfont{block title}{size=\normalsize}

% Badan blok rata kiri (bukan rata kanan-kiri): teks slide lebih mudah dibaca
% dan baris terakhir yang pendek tidak lagi memicu "Underfull \hbox".
\addtobeamertemplate{block body begin}{\raggedright}{}
\addtobeamertemplate{block alerted body begin}{\raggedright}{}
\addtobeamertemplate{block example body begin}{\raggedright}{}

% Header: nama mata kuliah dan kode MK
\setbeamertemplate{headline}{%
    \begin{beamercolorbox}[wd=\paperwidth,ht=2.6ex,dp=1.2ex,leftskip=1em,rightskip=1em]{palette primary}%
        \usebeamerfont{headline}\CourseName\ \textperiodcentered\ \CourseCode%
        \hfill \CourseProgram%
    \end{beamercolorbox}%
}

% Footer: minggu dan nomor halaman
\setbeamertemplate{footline}{%
    \begin{beamercolorbox}[wd=\paperwidth,ht=2.4ex,dp=1.2ex,leftskip=1em,rightskip=1em]{palette secondary}%
        \usebeamerfont{footline}\WeekLabel\ \textperiodcentered\ \WeekTopic%
        \hfill \insertframenumber\,/\,\inserttotalframenumber%
    \end{beamercolorbox}%
}

% Halaman judul: judul, subjudul, lalu blok identitas mata kuliah
\setbeamertemplate{title page}{%
    \begin{centering}
        \vskip1.0em
        {\usebeamerfont{title}\usebeamercolor[fg]{title}\inserttitle\par}
        \vskip0.3em
        {\usebeamerfont{subtitle}\usebeamercolor[fg]{subtitle}\insertsubtitle\par}
        \vskip1.2em
        \textcolor{ubbblue}{\rule{0.72\textwidth}{0.8pt}}
        \vskip1.2em
        {\small
            \textbf{\CourseName}\ (\CourseCode)\par
            \vskip0.4em
            \CourseProgram\par
            \CourseFaculty\par
            \CourseInstitution\par
            \vskip0.6em
            \CourseSemester\ \textperiodcentered\ \WeekLabel\par
            \vskip0.8em
            Dosen Pengampu: \CourseLecturer\par
        }
        \vskip1.0em
        \textcolor{ubbblue}{\rule{0.72\textwidth}{0.4pt}}
    \end{centering}
}

% =====================================================================
%  PERINTAH BANTU
% =====================================================================
\newcommand{\R}{\mathbb{R}}
\newcommand{\N}{\mathbb{N}}
\newcommand{\Z}{\mathbb{Z}}
\newcommand{\code}[1]{\texttt{#1}}

% Rujukan gambar dari buku: \bookfig[lebar]{NN}{namaberkas.png}
% Jalur relatif terhadap folder ppt/ (dikompilasi dari sana).
\newcommand{\bookfig}[3][0.80\linewidth]{%
    \includegraphics[width=#1]{../buku_ajar/chapters/#2/figures/#3}%
}

% Keterangan sumber gambar, sepadan dengan \sumber di buku
\newcommand{\figsumber}[1]{%
    \par\vspace{1pt}{\tiny\itshape Sumber: #1\par}%
}

% Penanda blok materi yang sering dipakai di slide
\newcommand{\slidejudul}[1]{{\large\bfseries\textcolor{ubbblue}{#1}}\par\vspace{2pt}}
\newcommand{\kunci}[1]{\textcolor{ubbred}{\textbf{#1}}}

%      \title{...} \subtitle{...} \setWeekInfo{Minggu n}{Topik}
%
%  Berkas minggu hanya berisi \setWeekInfo dan lingkungan frame.
%  Kompilasi dijalankan DARI DALAM folder ppt/ (lihat compile-ppt.bat),
%  karena gambar buku dirujuk dengan jalur relatif ../buku_ajar/.
% =====================================================================

% ---------------------------------------------------------------- bahasa
\usepackage[utf8]{inputenc}
\usepackage[T1]{fontenc}
\usepackage[indonesian]{babel}
\usepackage{lmodern}
\usepackage{microtype}

% ------------------------------------------------------------- matematika
\usepackage{amsmath,amssymb,mathtools}

% ----------------------------------------------------------------- tabel
\usepackage{booktabs}
\usepackage{array}
\usepackage{tabularx}

% --------------------------------------------------------------- grafik
\usepackage{tikz}
\usetikzlibrary{arrows.meta,positioning,shapes.geometric,calc,fit,backgrounds}

% ---------------------------------------------------------- kode program
\usepackage{listings}
\lstdefinestyle{slidecode}{
    basicstyle=\ttfamily\scriptsize,
    keywordstyle=\color{ubbblue}\bfseries,
    commentstyle=\color{ubbgreen}\itshape,
    stringstyle=\color{ubbred},
    numbers=left,
    numberstyle=\tiny\color{gray},
    numbersep=6pt,
    frame=single,
    rulecolor=\color{ubbblue!40},
    backgroundcolor=\color{ubbblue!3},
    breaklines=true,
    showstringspaces=false,
    tabsize=4,
    columns=fullflexible,
    keepspaces=true,
}
\lstset{style=slidecode}

% =====================================================================
%  WARNA INSTITUSI
% =====================================================================
\definecolor{ubbblue}{RGB}{12, 62, 120}      % warna utama
\definecolor{ubbteal}{RGB}{0, 118, 110}      % aksen kedua
\definecolor{ubbgreen}{RGB}{22, 122, 74}     % untuk blok contoh
\definecolor{ubbred}{RGB}{168, 42, 42}       % untuk blok peringatan
\definecolor{ubbgray}{RGB}{70, 76, 84}       % teks sekunder

% =====================================================================
%  IDENTITAS MATA KULIAH
% =====================================================================
\newcommand{\CourseName}{Pengantar Kriptografi}
\newcommand{\CourseCode}{MPFA014}
\newcommand{\CourseSemester}{Semester Ganjil}
\newcommand{\CourseProgram}{Program Studi Matematika}
\newcommand{\CourseFaculty}{Fakultas Matematika dan Ilmu Pengetahuan Alam}
\newcommand{\CourseInstitution}{Universitas Bale Bandung}
\newcommand{\CourseLecturer}{Cecep Suwanda, S.Si., M.Kom.}

% Label minggu, disetel tiap berkas lewat \setWeekInfo
\newcommand{\WeekLabel}{}
\newcommand{\WeekTopic}{}
\newcommand{\setWeekInfo}[2]{%
    \renewcommand{\WeekLabel}{#1}%
    \renewcommand{\WeekTopic}{#2}%
}

% =====================================================================
%  TEMA: header, footer, warna blok
% =====================================================================
\usetheme{default}
\setbeamertemplate{navigation symbols}{}
\setbeamertemplate{blocks}[rounded][shadow=false]

\setbeamercolor{palette primary}{bg=ubbblue, fg=white}
\setbeamercolor{palette secondary}{bg=ubbgray, fg=white}
\setbeamercolor{palette tertiary}{bg=ubbteal, fg=white}

\setbeamercolor{structure}{fg=ubbblue}
\setbeamercolor{normal text}{fg=black!92}
\setbeamercolor{alerted text}{fg=ubbred}

\setbeamercolor{block title}{bg=ubbblue, fg=white}
\setbeamercolor{block body}{bg=ubbblue!6, fg=black!92}
\setbeamercolor{block title alerted}{bg=ubbred, fg=white}
\setbeamercolor{block body alerted}{bg=ubbred!6, fg=black!92}
\setbeamercolor{block title example}{bg=ubbgreen, fg=white}
\setbeamercolor{block body example}{bg=ubbgreen!6, fg=black!92}

\setbeamercolor{itemize item}{fg=ubbblue}
\setbeamercolor{itemize subitem}{fg=ubbteal}
\setbeamercolor{enumerate item}{fg=ubbblue}

\setbeamerfont{headline}{size=\scriptsize}
\setbeamerfont{footline}{size=\scriptsize}
\setbeamerfont{block title}{size=\normalsize}

% Badan blok rata kiri (bukan rata kanan-kiri): teks slide lebih mudah dibaca
% dan baris terakhir yang pendek tidak lagi memicu "Underfull \hbox".
\addtobeamertemplate{block body begin}{\raggedright}{}
\addtobeamertemplate{block alerted body begin}{\raggedright}{}
\addtobeamertemplate{block example body begin}{\raggedright}{}

% Header: nama mata kuliah dan kode MK
\setbeamertemplate{headline}{%
    \begin{beamercolorbox}[wd=\paperwidth,ht=2.6ex,dp=1.2ex,leftskip=1em,rightskip=1em]{palette primary}%
        \usebeamerfont{headline}\CourseName\ \textperiodcentered\ \CourseCode%
        \hfill \CourseProgram%
    \end{beamercolorbox}%
}

% Footer: minggu dan nomor halaman
\setbeamertemplate{footline}{%
    \begin{beamercolorbox}[wd=\paperwidth,ht=2.4ex,dp=1.2ex,leftskip=1em,rightskip=1em]{palette secondary}%
        \usebeamerfont{footline}\WeekLabel\ \textperiodcentered\ \WeekTopic%
        \hfill \insertframenumber\,/\,\inserttotalframenumber%
    \end{beamercolorbox}%
}

% Halaman judul: judul, subjudul, lalu blok identitas mata kuliah
\setbeamertemplate{title page}{%
    \begin{centering}
        \vskip1.0em
        {\usebeamerfont{title}\usebeamercolor[fg]{title}\inserttitle\par}
        \vskip0.3em
        {\usebeamerfont{subtitle}\usebeamercolor[fg]{subtitle}\insertsubtitle\par}
        \vskip1.2em
        \textcolor{ubbblue}{\rule{0.72\textwidth}{0.8pt}}
        \vskip1.2em
        {\small
            \textbf{\CourseName}\ (\CourseCode)\par
            \vskip0.4em
            \CourseProgram\par
            \CourseFaculty\par
            \CourseInstitution\par
            \vskip0.6em
            \CourseSemester\ \textperiodcentered\ \WeekLabel\par
            \vskip0.8em
            Dosen Pengampu: \CourseLecturer\par
        }
        \vskip1.0em
        \textcolor{ubbblue}{\rule{0.72\textwidth}{0.4pt}}
    \end{centering}
}

% =====================================================================
%  PERINTAH BANTU
% =====================================================================
\newcommand{\R}{\mathbb{R}}
\newcommand{\N}{\mathbb{N}}
\newcommand{\Z}{\mathbb{Z}}
\newcommand{\code}[1]{\texttt{#1}}

% Rujukan gambar dari buku: \bookfig[lebar]{NN}{namaberkas.png}
% Jalur relatif terhadap folder ppt/ (dikompilasi dari sana).
\newcommand{\bookfig}[3][0.80\linewidth]{%
    \includegraphics[width=#1]{../buku_ajar/chapters/#2/figures/#3}%
}

% Keterangan sumber gambar, sepadan dengan \sumber di buku
\newcommand{\figsumber}[1]{%
    \par\vspace{1pt}{\tiny\itshape Sumber: #1\par}%
}

% Penanda blok materi yang sering dipakai di slide
\newcommand{\slidejudul}[1]{{\large\bfseries\textcolor{ubbblue}{#1}}\par\vspace{2pt}}
\newcommand{\kunci}[1]{\textcolor{ubbred}{\textbf{#1}}}

%      \title{...} \subtitle{...} \setWeekInfo{Minggu n}{Topik}
%
%  Berkas minggu hanya berisi \setWeekInfo dan lingkungan frame.
%  Kompilasi dijalankan DARI DALAM folder ppt/ (lihat compile-ppt.bat),
%  karena gambar buku dirujuk dengan jalur relatif ../buku_ajar/.
% =====================================================================

% ---------------------------------------------------------------- bahasa
\usepackage[utf8]{inputenc}
\usepackage[T1]{fontenc}
\usepackage[indonesian]{babel}
\usepackage{lmodern}
\usepackage{microtype}

% ------------------------------------------------------------- matematika
\usepackage{amsmath,amssymb,mathtools}

% ----------------------------------------------------------------- tabel
\usepackage{booktabs}
\usepackage{array}
\usepackage{tabularx}

% --------------------------------------------------------------- grafik
\usepackage{tikz}
\usetikzlibrary{arrows.meta,positioning,shapes.geometric,calc,fit,backgrounds}

% ---------------------------------------------------------- kode program
\usepackage{listings}
\lstdefinestyle{slidecode}{
    basicstyle=\ttfamily\scriptsize,
    keywordstyle=\color{ubbblue}\bfseries,
    commentstyle=\color{ubbgreen}\itshape,
    stringstyle=\color{ubbred},
    numbers=left,
    numberstyle=\tiny\color{gray},
    numbersep=6pt,
    frame=single,
    rulecolor=\color{ubbblue!40},
    backgroundcolor=\color{ubbblue!3},
    breaklines=true,
    showstringspaces=false,
    tabsize=4,
    columns=fullflexible,
    keepspaces=true,
}
\lstset{style=slidecode}

% =====================================================================
%  WARNA INSTITUSI
% =====================================================================
\definecolor{ubbblue}{RGB}{12, 62, 120}      % warna utama
\definecolor{ubbteal}{RGB}{0, 118, 110}      % aksen kedua
\definecolor{ubbgreen}{RGB}{22, 122, 74}     % untuk blok contoh
\definecolor{ubbred}{RGB}{168, 42, 42}       % untuk blok peringatan
\definecolor{ubbgray}{RGB}{70, 76, 84}       % teks sekunder

% =====================================================================
%  IDENTITAS MATA KULIAH
% =====================================================================
\newcommand{\CourseName}{Pengantar Kriptografi}
\newcommand{\CourseCode}{MPFA014}
\newcommand{\CourseSemester}{Semester Ganjil}
\newcommand{\CourseProgram}{Program Studi Matematika}
\newcommand{\CourseFaculty}{Fakultas Matematika dan Ilmu Pengetahuan Alam}
\newcommand{\CourseInstitution}{Universitas Bale Bandung}
\newcommand{\CourseLecturer}{Cecep Suwanda, S.Si., M.Kom.}

% Label minggu, disetel tiap berkas lewat \setWeekInfo
\newcommand{\WeekLabel}{}
\newcommand{\WeekTopic}{}
\newcommand{\setWeekInfo}[2]{%
    \renewcommand{\WeekLabel}{#1}%
    \renewcommand{\WeekTopic}{#2}%
}

% =====================================================================
%  TEMA: header, footer, warna blok
% =====================================================================
\usetheme{default}
\setbeamertemplate{navigation symbols}{}
\setbeamertemplate{blocks}[rounded][shadow=false]

\setbeamercolor{palette primary}{bg=ubbblue, fg=white}
\setbeamercolor{palette secondary}{bg=ubbgray, fg=white}
\setbeamercolor{palette tertiary}{bg=ubbteal, fg=white}

\setbeamercolor{structure}{fg=ubbblue}
\setbeamercolor{normal text}{fg=black!92}
\setbeamercolor{alerted text}{fg=ubbred}

\setbeamercolor{block title}{bg=ubbblue, fg=white}
\setbeamercolor{block body}{bg=ubbblue!6, fg=black!92}
\setbeamercolor{block title alerted}{bg=ubbred, fg=white}
\setbeamercolor{block body alerted}{bg=ubbred!6, fg=black!92}
\setbeamercolor{block title example}{bg=ubbgreen, fg=white}
\setbeamercolor{block body example}{bg=ubbgreen!6, fg=black!92}

\setbeamercolor{itemize item}{fg=ubbblue}
\setbeamercolor{itemize subitem}{fg=ubbteal}
\setbeamercolor{enumerate item}{fg=ubbblue}

\setbeamerfont{headline}{size=\scriptsize}
\setbeamerfont{footline}{size=\scriptsize}
\setbeamerfont{block title}{size=\normalsize}

% Badan blok rata kiri (bukan rata kanan-kiri): teks slide lebih mudah dibaca
% dan baris terakhir yang pendek tidak lagi memicu "Underfull \hbox".
\addtobeamertemplate{block body begin}{\raggedright}{}
\addtobeamertemplate{block alerted body begin}{\raggedright}{}
\addtobeamertemplate{block example body begin}{\raggedright}{}

% Header: nama mata kuliah dan kode MK
\setbeamertemplate{headline}{%
    \begin{beamercolorbox}[wd=\paperwidth,ht=2.6ex,dp=1.2ex,leftskip=1em,rightskip=1em]{palette primary}%
        \usebeamerfont{headline}\CourseName\ \textperiodcentered\ \CourseCode%
        \hfill \CourseProgram%
    \end{beamercolorbox}%
}

% Footer: minggu dan nomor halaman
\setbeamertemplate{footline}{%
    \begin{beamercolorbox}[wd=\paperwidth,ht=2.4ex,dp=1.2ex,leftskip=1em,rightskip=1em]{palette secondary}%
        \usebeamerfont{footline}\WeekLabel\ \textperiodcentered\ \WeekTopic%
        \hfill \insertframenumber\,/\,\inserttotalframenumber%
    \end{beamercolorbox}%
}

% Halaman judul: judul, subjudul, lalu blok identitas mata kuliah
\setbeamertemplate{title page}{%
    \begin{centering}
        \vskip1.0em
        {\usebeamerfont{title}\usebeamercolor[fg]{title}\inserttitle\par}
        \vskip0.3em
        {\usebeamerfont{subtitle}\usebeamercolor[fg]{subtitle}\insertsubtitle\par}
        \vskip1.2em
        \textcolor{ubbblue}{\rule{0.72\textwidth}{0.8pt}}
        \vskip1.2em
        {\small
            \textbf{\CourseName}\ (\CourseCode)\par
            \vskip0.4em
            \CourseProgram\par
            \CourseFaculty\par
            \CourseInstitution\par
            \vskip0.6em
            \CourseSemester\ \textperiodcentered\ \WeekLabel\par
            \vskip0.8em
            Dosen Pengampu: \CourseLecturer\par
        }
        \vskip1.0em
        \textcolor{ubbblue}{\rule{0.72\textwidth}{0.4pt}}
    \end{centering}
}

% =====================================================================
%  PERINTAH BANTU
% =====================================================================
\newcommand{\R}{\mathbb{R}}
\newcommand{\N}{\mathbb{N}}
\newcommand{\Z}{\mathbb{Z}}
\newcommand{\code}[1]{\texttt{#1}}

% Rujukan gambar dari buku: \bookfig[lebar]{NN}{namaberkas.png}
% Jalur relatif terhadap folder ppt/ (dikompilasi dari sana).
\newcommand{\bookfig}[3][0.80\linewidth]{%
    \includegraphics[width=#1]{../buku_ajar/chapters/#2/figures/#3}%
}

% Keterangan sumber gambar, sepadan dengan \sumber di buku
\newcommand{\figsumber}[1]{%
    \par\vspace{1pt}{\tiny\itshape Sumber: #1\par}%
}

% Penanda blok materi yang sering dipakai di slide
\newcommand{\slidejudul}[1]{{\large\bfseries\textcolor{ubbblue}{#1}}\par\vspace{2pt}}
\newcommand{\kunci}[1]{\textcolor{ubbred}{\textbf{#1}}}


\title{Cipher Klasik I: Substitusi dan Transposisi}
\subtitle{Dari Caesar dan Affine sampai Rail Fence dan super-enkripsi}
\setWeekInfo{Minggu 4}{Kriptografi klasik I}

\begin{document}

% ----------------------------------------------------------- halaman judul
\begin{frame}[plain]
    \titlepage
\end{frame}

% ------------------------------------------------------------------ peta
\begin{frame}{Peta Materi Pertemuan Ini}
    \begin{columns}[T,onlytextwidth]
        \begin{column}{0.54\textwidth}
            \slidejudul{Lima bagian}
            \begin{enumerate}
                \item Kriptografi klasik dan dua teknik dasar
                \item Cipher substitusi monoalfabetik dan Caesar
                \item Affine cipher
                \item Transposisi Rail Fence
                \item Super-enkripsi dan kelemahan substitusi
            \end{enumerate}
        \end{column}
        \begin{column}{0.42\textwidth}
            \begin{block}{Referensi}
                \footnotesize
                Bab 4 buku ajar; Munir, \textit{IF4020 Kriptografi},
                ITB, 2026 (modul 03 dan 04)
            \end{block}
            \begin{exampleblock}{Benang merah}
                \footnotesize
                Substitusi menyembunyikan identitas huruf; transposisi
                menyembunyikan urutannya.
            \end{exampleblock}
        \end{column}
    \end{columns}
\end{frame}

% ---------------------------------------------------- capaian pembelajaran
\begin{frame}{Capaian Pembelajaran}
    \begin{block}{Sub-CPMK yang dibebankan}
        \footnotesize
        \textbf{Sub-CPMK 1.1} --- Mengenal cipher klasik, yaitu substitusi
        (Caesar, monoalfabetik, Affine), transposisi (Rail Fence),
        Vigen\`{e}re, Playfair, Hill, dan Enigma, beserta gagasan
        kriptanalisisnya (analisis frekuensi dan metode Kasiski).
    \end{block}
    \vspace{0.15cm}
    \begin{exampleblock}{Pada pertemuan ini, mahasiswa mampu}
        \begin{itemize}
            \item menjelaskan perbedaan cipher substitusi dan transposisi;
            \item mengenal cipher monoalfabetik dan Caesar beserta
                  kriptanalisisnya;
            \item menjelaskan rumus dan syarat kunci Affine;
            \item mengidentifikasi pola zigzag pada transposisi Rail Fence;
            \item menjelaskan super-enkripsi sebagai gabungan dua teknik.
        \end{itemize}
    \end{exampleblock}
\end{frame}

% =====================================================================
\section{Kriptografi Klasik dan Dua Teknik}
% =====================================================================

\begin{frame}{Kriptografi Klasik: Ciri dan Alasannya}
    \begin{itemize}
        \item Bekerja pada huruf alfabet, bukan pada bit atau blok data.
        \item Alat utamanya pena dan kertas; belum ada komputer digital.
        \item Termasuk kriptografi kunci-simetri: pengirim dan penerima
              memakai kunci yang sama.
        \item Usianya ribuan tahun, tetapi gagasannya masih hidup di cipher
              modern.
    \end{itemize}
    \vspace{0.15cm}
    \begin{block}{Tiga alasan mempelajarinya}
        \begin{enumerate}
            \item Memahami konsep dasar kriptografi.
            \item Melihat akar algoritma kriptografi modern.
            \item Mengenali kelemahan yang membuat suatu cipher mudah
                  dipecahkan.
        \end{enumerate}
    \end{block}
\end{frame}

\begin{frame}{Dua Teknik Dasar}
    \begin{columns}[T,onlytextwidth]
        \begin{column}{0.49\textwidth}
            \begin{block}{Substitusi}
                \footnotesize
                Mengganti \emph{identitas} huruf; urutan huruf tidak
                berubah.\\[3pt]
                \code{MENGANTUK} $\longrightarrow$ \code{CQBSIBONW}
            \end{block}
        \end{column}
        \begin{column}{0.49\textwidth}
            \begin{exampleblock}{Transposisi}
                \footnotesize
                Memindahkan \emph{posisi} huruf; identitas huruf tidak
                berubah.\\[3pt]
                \code{MENGANTUK} $\longrightarrow$ \code{TNEAKMNGU}
            \end{exampleblock}
        \end{column}
    \end{columns}
    \vspace{0.3cm}
    \begin{alertblock}{Gabungan keduanya}
        Substitusi lalu transposisi, atau sebaliknya, disebut
        \textbf{super-enkripsi} atau \textit{product cipher}.
    \end{alertblock}
\end{frame}

\begin{frame}{Peta Kriptologi}
    \begin{columns}[T,onlytextwidth]
        \begin{column}{0.55\textwidth}
            \begin{itemize}
                \item Kriptografi klasik bercabang menjadi cipher substitusi
                      dan cipher transposisi.
                \item Setiap cabang dapat dipakai sendiri-sendiri.
                \item Keduanya juga dapat dirangkai menjadi super-enkripsi.
            \end{itemize}
            \vspace{0.1cm}
            \begin{block}{Fokus minggu ini}
                \footnotesize
                Monoalfabetik, Caesar, Affine, Rail Fence, dan super-enkripsi.
            \end{block}
        \end{column}
        \begin{column}{0.42\textwidth}
            \centering
            \bookfig[0.98\linewidth]{04}{kriptologi-hierarki.png}
            \figsumber{Munir, \textit{IF4020 Kriptografi}, ITB.}
        \end{column}
    \end{columns}
\end{frame}

% =====================================================================
\section{Cipher Substitusi Monoalfabetik}
% =====================================================================

\begin{frame}{Cipher Abjad-Tunggal}
    Pada cipher abjad-tunggal, setiap huruf plainteks dipetakan ke tepat satu
    huruf cipherteks, dan pemetaan itu tetap untuk seluruh pesan.
    \[
        26! = 403.291.461.126.605.635.584.000.000
    \]
    \begin{itemize}
        \item Pemetaan harus satu-satu: 26 huruf menempati 26 posisi berbeda.
        \item Banyaknya tabel substitusi yang mungkin adalah $26!$.
        \item Tabel dapat dibentuk dari permutasi sembarang atau dari kalimat
              kunci.
    \end{itemize}
    \begin{alertblock}{Ruang kunci besar bukan jaminan}
        \footnotesize
        Struktur statistik bahasa tetap lolos ke cipherteks, sehingga cipher
        ini dapat dipecahkan tanpa mencoba seluruh kunci.
    \end{alertblock}
\end{frame}

\begin{frame}{Tabel Substitusi dari Kalimat Kunci}
    \footnotesize
    Kalimat kunci: \code{di bawah sinar bulan purnama hati resah jadi senang}.
    Huruf berulang dibuang sehingga tersisa 17 huruf:
    \code{DIBAWHSNRULPMTEJG}. Sisa alfabet disambung menurut abjad:
    \code{CFKOQVXYZ}.
    \vspace{0.1cm}
    \begin{center}
    \scriptsize
    \setlength{\tabcolsep}{3pt}
    \begin{tabular}{@{}l*{13}{c}@{}}
        \toprule
        Plainteks  & A&B&C&D&E&F&G&H&I&J&K&L&M \\
        Cipherteks & D&I&B&A&W&H&S&N&R&U&L&P&M \\
        \midrule
        Plainteks  & N&O&P&Q&R&S&T&U&V&W&X&Y&Z \\
        Cipherteks & T&E&J&G&C&F&K&O&Q&V&X&Y&Z \\
        \bottomrule
    \end{tabular}
    \end{center}
    \vspace{0.05cm}
    \begin{block}{Yang perlu dicatat}
        \footnotesize
        Kalimat kunci hanya memudahkan menghafal tabel; ia tidak menambah
        banyaknya kemungkinan kunci.
    \end{block}
\end{frame}

\begin{frame}{Caesar Cipher: Pergeseran Tetap}
    \begin{columns}[T,onlytextwidth]
        \begin{column}{0.54\textwidth}
            Caesar cipher menggeser setiap huruf sejauh $k$ posisi. Versi
            historis memakai $k=3$.
            \begin{itemize}
                \item \code{A} $\rightarrow$ \code{D},
                      \code{B} $\rightarrow$ \code{E}, dan seterusnya.
                \item \code{X}, \code{Y}, \code{Z} melingkar ke
                      \code{A}, \code{B}, \code{C}.
            \end{itemize}
            Dengan pemetaan $A=0,\dots,Z=25$:
            \[
                c=(p+k)\bmod 26, \qquad p=(c-k)\bmod 26.
            \]
            \footnotesize
            Untuk alfabet 256 karakter ASCII, modulusnya menjadi 256.
        \end{column}
        \begin{column}{0.43\textwidth}
            \centering
            \bookfig[0.98\linewidth]{04}{caesar-pergeseran.png}
            \figsumber{Munir, \textit{IF4020 Kriptografi}, ITB.}
        \end{column}
    \end{columns}
\end{frame}

\begin{frame}{Alat Bantu: Cakram Caesar}
    \begin{columns}[T,onlytextwidth]
        \begin{column}{0.58\textwidth}
            \begin{itemize}
                \item Cakram Caesar memakai dua lingkar alfabet yang dapat
                      diputar.
                \item Memutar lingkar luar sejauh $k$ posisi sama dengan
                      menambahkan $k$ pada setiap huruf.
                \item Hasilnya identik dengan rumus $(p+k)\bmod 26$.
                \item Alat ini membantu penghitungan manual, bukan menambah
                      keamanan.
            \end{itemize}
        \end{column}
        \begin{column}{0.40\textwidth}
            \centering
            \bookfig[0.88\linewidth]{04}{cakram-caesar.png}
            \figsumber{Munir, \textit{IF4020 Kriptografi}, ITB.}
        \end{column}
    \end{columns}
\end{frame}

\begin{frame}{Contoh Caesar: Enkripsi dengan $k=3$}
    \footnotesize
    Plainteks: \code{temui saya di jembatan merah nanti malam}. Spasi
    dibuang, lalu setiap huruf ditambah 3 modulo 26.
    \vspace{0.1cm}
    \begin{columns}[T,onlytextwidth]
        \begin{column}{0.50\textwidth}
            \begin{center}
            \begin{tabular}{@{}c c c c@{}}
                \toprule
                Huruf & $p$ & $(p+3)\bmod 26$ & Hasil \\
                \midrule
                \code{t} & 19 & 22 & \code{W} \\
                \code{e} & 4  & 7  & \code{H} \\
                \code{m} & 12 & 15 & \code{P} \\
                \code{u} & 20 & 23 & \code{X} \\
                \code{i} & 8  & 11 & \code{L} \\
                \bottomrule
            \end{tabular}
            \end{center}
        \end{column}
        \begin{column}{0.48\textwidth}
            \begin{exampleblock}{Cipherteks lengkap}
                \code{WHPXLVDBDGLMHPEDWDQPHUDKQDQWLPDODP}\\[3pt]
                Dikelompokkan empat huruf:\\
                \code{WHPX LVDB DGLM HPED WDQP}\\
                \code{HUDK QDQW LPDO DP}
            \end{exampleblock}
        \end{column}
    \end{columns}
    \vspace{0.05cm}
    {\scriptsize Lima huruf pertama \code{temui} menjadi \code{WHPXL}; 34
    huruf plainteks menghasilkan 34 huruf cipherteks.}
\end{frame}

\begin{frame}{Dekripsi Caesar dan Konvensi Penulisan}
    \begin{columns}[T,onlytextwidth]
        \begin{column}{0.52\textwidth}
            Dekripsi memakai $p=(c-k)\bmod 26$:
            \begin{itemize}
                \item \code{W} $=22 \rightarrow (22-3)\bmod 26=19=\code{t}$
                \item \code{H} $=7 \rightarrow (7-3)\bmod 26=4=\code{e}$
                \item \code{P} $=15 \rightarrow (15-3)\bmod 26=12=\code{m}$
                \item \code{X} $=23 \rightarrow (23-3)\bmod 26=20=\code{u}$
                \item \code{L} $=11 \rightarrow (11-3)\bmod 26=8=\code{i}$
            \end{itemize}
        \end{column}
        \begin{column}{0.46\textwidth}
            \begin{block}{Konvensi cipherteks}
                \footnotesize
                Spasi dibuang, lalu cipherteks dikelompokkan menjadi blok
                empat huruf: \code{WHPX LVDB DGLM \dots}
            \end{block}
            \begin{alertblock}{Tujuannya}
                \footnotesize
                Menghapus petunjuk panjang kata dan membuat cipherteks lebih
                sulit dibaca.
            \end{alertblock}
        \end{column}
    \end{columns}
\end{frame}

\begin{frame}{Kriptanalisis Caesar: Coba Semua Kunci}
    \footnotesize
    Hanya ada 26 kemungkinan kunci. Penyerang mendekripsi cipherteks
    \code{XMZVH} dengan setiap nilai $k$ lalu memeriksa hasilnya.
    \vspace{0.1cm}
    \begin{center}
    \scriptsize
    \begin{tabular}{@{}l l @{\hspace{14pt}} l l@{}}
        \toprule
        $k$ & Dekripsi & $k$ & Dekripsi \\
        \midrule
        0  & \code{XMZVH} & 13 & \code{KZMIU} \\
        1  & \code{WLYUG} & 14 & \code{JYLHT} \\
        2  & \code{VKXTF} & 15 & \code{IXKGS} \\
        3  & \code{UJWSE} & 16 & \code{HWJFR} \\
        4  & \code{TIVRD} & 17 & \code{GVIEQ} \\
        5  & \code{SHUQC} & 18 & \code{FUHDP} \\
        6  & \code{RGTPB} & 19 & \code{ETGCO} \\
        7  & \code{QFSOA} & 20 & \code{DSFBN} \\
        8  & \code{PERNZ} & 21 & \textbf{\code{CREAM}} \\
        9  & \code{ODQMY} & 22 & \code{BQDZL} \\
        10 & \code{NCPLX} & 23 & \code{APCYK} \\
        11 & \code{MBOKW} & 24 & \code{ZOBXJ} \\
        12 & \code{LANJV} & 25 & \code{YNAWI} \\
        \bottomrule
    \end{tabular}
    \end{center}
    \vspace{0.02cm}
    {\scriptsize Hanya $k=21$ yang memberi kata bermakna, yaitu \code{CREAM};
    jadi kunci yang benar adalah 21.}
\end{frame}

\begin{frame}{Bila Dua Kunci Sama-Sama Bermakna}
    \begin{itemize}
        \item Cipherteks \code{HSPPW} pada $k=4$ menghasilkan \code{DOLLS}.
        \item Pada $k=11$ cipherteks yang sama menghasilkan \code{WHEEL}.
        \item Keduanya kata bermakna, sehingga pencarian menyeluruh saja tidak
              cukup.
        \item Konteks pesan atau potongan cipherteks lain dipakai untuk memilih
              kunci.
    \end{itemize}
    \vspace{0.15cm}
    \begin{alertblock}{Pelajaran}
        \footnotesize
        Caesar tetap rapuh karena ruang kuncinya hanya 26; ambiguitas hanya
        mempersulit pemeriksaan akhir, bukan menyelamatkannya.
    \end{alertblock}
\end{frame}

\begin{frame}{ROT13: Caesar dengan $k=13$}
    \begin{columns}[T,onlytextwidth]
        \begin{column}{0.55\textwidth}
            \begin{itemize}
                \item ROT13 adalah Caesar cipher dengan pergeseran 13.
                \item Karena $13+13=26$, penerapan dua kali mengembalikan
                      pesan semula.
                \item $\mathrm{ROT13}(\code{ROTATE})=\code{EBGNGR}$.
                \item Enkripsi dan dekripsi memakai operasi yang sama.
            \end{itemize}
            \[
                \mathrm{ROT13}(\mathrm{ROT13}(P))=P.
            \]
            \begin{alertblock}{Bukan untuk kerahasiaan}
                \footnotesize
                ROT13 dipakai untuk menyembunyikan jawaban teka-teki atau
                lelucon, bukan untuk mengamankan pesan.
            \end{alertblock}
        \end{column}
        \begin{column}{0.43\textwidth}
            \centering
            \bookfig[0.98\linewidth]{04}{rot13.png}
            \figsumber{Munir, \textit{IF4020 Kriptografi}, ITB.}
        \end{column}
    \end{columns}
\end{frame}

% =====================================================================
\section{Affine Cipher}
% =====================================================================

\begin{frame}{Dari Caesar ke Affine}
    Caesar hanya menggeser setiap huruf. Affine menambahkan satu operasi:
    mengalikan nilai huruf dengan $m$ sebelum menggeser sejauh $b$.
    \begin{block}{Rumus}
        \[
            C \equiv mP+b \pmod{26}, \qquad
            P \equiv m^{-1}(C-b) \pmod{26}.
        \]
    \end{block}
    \begin{itemize}
        \item Kuncinya adalah pasangan $(m,b)$.
        \item $m^{-1}$ adalah invers $m$ modulo 26, yaitu
              $m\cdot m^{-1}\equiv 1\pmod{26}$.
        \item Caesar adalah kasus khusus Affine dengan $m=1$.
    \end{itemize}
\end{frame}

\begin{frame}{Syarat Kunci Affine}
    \begin{itemize}
        \item Ukuran alfabet $n=26$.
        \item $m$ harus relatif prima dengan 26 agar pemetaan satu-satu dan
              inversnya ada.
        \item Nilai $m$ yang sah: $1,3,5,7,9,11,15,17,19,21,23,25$
              (12 nilai).
        \item Ada 25 pilihan untuk $b$.
        \item Ruang kunci: $12\times 25=300$ kemungkinan.
    \end{itemize}
    \vspace{0.1cm}
    \begin{alertblock}{Bandingkan dengan cipher abjad-tunggal}
        \footnotesize
        $26!$ tabel menjadi hanya 300 kunci; Affine jauh lebih mudah dicoba
        seluruhnya.
    \end{alertblock}
\end{frame}

\begin{frame}{Contoh Affine: Enkripsi \code{kripto}}
    \footnotesize
    Kunci $m=7$ dan $b=10$, sehingga $C\equiv 7P+10\pmod{26}$.
    \vspace{0.1cm}
    \begin{center}
    \begin{tabular}{@{}c c c c c@{}}
        \toprule
        Plainteks & $P$ & $7P+10$ & $(7P+10)\bmod 26$ & Cipherteks \\
        \midrule
        \code{k} & 10 & 80  & 2  & \code{C} \\
        \code{r} & 17 & 129 & 25 & \code{Z} \\
        \code{i} & 8  & 66  & 14 & \code{O} \\
        \code{p} & 15 & 115 & 11 & \code{L} \\
        \code{t} & 19 & 143 & 13 & \code{N} \\
        \code{o} & 14 & 108 & 4  & \code{E} \\
        \bottomrule
    \end{tabular}
    \end{center}
    \vspace{0.05cm}
    \begin{exampleblock}{Hasil}
        Cipherteks \code{kripto} adalah \code{CZOLNE}. Setiap baris memuat
        operasi perkalian, penambahan, dan sisa bagi.
    \end{exampleblock}
\end{frame}

\begin{frame}{Contoh Affine: Dekripsi dengan Invers}
    \footnotesize
    Karena $7\cdot 15=105\equiv 1\pmod{26}$, inversnya $7^{-1}=15$.
    Dekripsi memakai $P\equiv 15(C-10)\pmod{26}$.
    \vspace{0.1cm}
    \begin{center}
    \begin{tabular}{@{}c c c c c@{}}
        \toprule
        Cipherteks & $C$ & $15(C-10)$ & $\bmod 26$ & Plainteks \\
        \midrule
        \code{C} & 2  & $-120$ & 10 & \code{k} \\
        \code{Z} & 25 & 225    & 17 & \code{r} \\
        \code{O} & 14 & 60     & 8  & \code{i} \\
        \code{L} & 11 & 15     & 15 & \code{p} \\
        \code{N} & 13 & 45     & 19 & \code{t} \\
        \code{E} & 4  & $-90$  & 14 & \code{o} \\
        \bottomrule
    \end{tabular}
    \end{center}
    \vspace{0.05cm}
    {\scriptsize Hasilnya kembali \code{kripto}. Perhatikan bahwa
    $(-120)\bmod 26=10$ dan $(-90)\bmod 26=14$; sisa bagi selalu diambil
    nonnegatif.}
\end{frame}

\begin{frame}{Kriptanalisis Affine: Known-Plaintext}
    \footnotesize
    Misalkan kriptanalis mengetahui dua pasangan: cipherteks \code{C} berasal
    dari plainteks \code{K}, dan \code{E} dari \code{O}. Nilainya
    $C=2$, $K=10$, $E=4$, $O=14$.
    \begin{block}{Dua kekongruenan}
        \[
            2\equiv 10m+b\pmod{26}, \qquad
            4\equiv 14m+b\pmod{26}.
        \]
    \end{block}
    \begin{itemize}
        \item Kurangkan persamaan kedua dengan pertama:
              $2\equiv 4m\pmod{26}$.
        \item Solusi yang dipakai: $m=7$.
        \item Substitusi $m=7$: $2\equiv 70+b\pmod{26}$, sehingga $b=10$.
        \item Kunci yang diperoleh: $(m,b)=(7,10)$.
    \end{itemize}
    \begin{alertblock}{Mengapa berhasil}
        \footnotesize
        Dua pasangan huruf memberi dua persamaan untuk dua bilangan yang tidak
        diketahui.
    \end{alertblock}
\end{frame}

\begin{frame}{Kriptanalisis Affine: Pencarian Menyeluruh}
    Bila pasangan plainteks-cipherteks tidak tersedia, penyerang masih dapat
    mencoba seluruh kunci.
    \begin{itemize}
        \item 12 nilai $m$ yang sah dikalikan 25 nilai $b$ memberi 300 kunci.
        \item Setiap kunci menghasilkan satu kandidat plainteks.
        \item Kandidat yang bermakna menunjukkan kunci yang benar.
        \item Komputer mencoba 300 kunci dalam waktu yang sangat singkat.
    \end{itemize}
    \vspace{0.15cm}
    \begin{block}{Pelajaran}
        \footnotesize
        Keamanan cipher tidak ditentukan oleh kerumitan rumusnya, melainkan
        oleh besar ruang kunci dan ketahanan strukturnya.
    \end{block}
\end{frame}

\begin{frame}{Varian: Affine atas Blok Huruf}
    \footnotesize
    Untuk memperbesar ruang kunci, huruf dienkripsi per blok, bukan satu per
    satu.
    \begin{itemize}
        \item Pesan \code{kriptografi} dipecah menjadi \code{krip},
              \code{togr}, \code{afi}.
        \item Setiap huruf menjadi dua angka:
              \code{krip} $\rightarrow 10\,17\,08\,15=10170815$.
        \item Blok terbesar \code{ZZZZ} bernilai $25252525$, dipakai sebagai
              modulus.
    \end{itemize}
    \begin{block}{Fungsi enkripsi}
        \[
            C\equiv 21\,035\,433\,P+23\,210\,025 \pmod{25\,252\,525}.
        \]
    \end{block}
    \begin{exampleblock}{Contoh dan dekripsi}
        $P=10170815$ menghasilkan $C=22\,837\,395$. Dekripsi memakai invers
        $11\,837\,547$.
    \end{exampleblock}
\end{frame}

% =====================================================================
\section{Transposisi Rail Fence}
% =====================================================================

\begin{frame}{Cipher Transposisi: Mengubah Posisi}
    Cipher transposisi tidak mengganti huruf, tetapi memindahkan posisinya.
    \begin{columns}[T,onlytextwidth]
        \begin{column}{0.48\textwidth}
            \begin{block}{Yang tetap}
                \footnotesize
                Panjang pesan dan sebaran frekuensi huruf tunggal tidak
                berubah.
            \end{block}
        \end{column}
        \begin{column}{0.48\textwidth}
            \begin{alertblock}{Yang berubah}
                \footnotesize
                Urutan huruf, sehingga pola bigram dan trigram ikut terganggu.
            \end{alertblock}
        \end{column}
    \end{columns}
    \vspace{0.15cm}
    \begin{itemize}
        \item Nama lain: cipher permutasi atau \textit{scrambling}.
        \item Buku membahas transposisi kolom, Rail Fence, dan transposisi
              blok.
        \item Pertemuan ini berfokus pada Rail Fence.
    \end{itemize}
\end{frame}

\begin{frame}{Rail Fence: Pola Zigzag}
    \begin{columns}[T,onlytextwidth]
        \begin{column}{0.49\textwidth}
            \begin{itemize}
                \item Plainteks ditulis menuruni dan menaiki beberapa baris.
                \item Setiap baris disebut \textit{rail} atau pagar.
                \item Untuk $k=3$ baris, pola nomor barisnya
                      $1,2,3,2,1,2,3,\dots$
            \end{itemize}
        \end{column}
        \begin{column}{0.49\textwidth}
            \begin{itemize}
                \item Cipherteks dibaca baris demi baris dari atas.
                \item Kuncinya adalah banyaknya baris $k$.
            \end{itemize}
        \end{column}
    \end{columns}
    \vspace{0.1cm}
    \begin{center}
        \bookfig[0.98\linewidth]{04}{rail-fence-3baris.png}
        \figsumber{Munir, \textit{IF4020 Kriptografi}, ITB.}
    \end{center}
\end{frame}

\begin{frame}{Contoh Rail Fence: $k=3$}
    \footnotesize
    Plainteks \code{CRYPTOGRAPHY AND DATA SECURITY}; spasi dibuang sehingga
    tersisa 27 huruf.
    \vspace{0.05cm}
    \begin{center}
    \begin{tabular}{@{}l l@{}}
        \toprule
        Baris & Isi \\
        \midrule
        1 & \code{CTAAAEI} \\
        2 & \code{RPORPYNDTSCRT} \\
        3 & \code{YGHDAUY} \\
        \bottomrule
    \end{tabular}
    \end{center}
    \vspace{0.02cm}
    {\scriptsize Pola penempatan: \code{C} ke baris 1, \code{R} ke baris 2,
    \code{Y} ke baris 3, \code{P} kembali ke baris 2, \code{T} ke baris 1,
    dan seterusnya.}
    \vspace{0.08cm}
    \begin{exampleblock}{Cipherteks}
        \code{CTAAAEIRPORPYNDTSCRTYGHDAUY}
    \end{exampleblock}
\end{frame}

\begin{frame}{Kelemahan Rail Fence}
    \begin{itemize}
        \item Kuncinya hanya bilangan $k$, yaitu banyaknya baris.
        \item Untuk 27 huruf, penyerang cukup mencoba $k=2$ sampai $26$.
        \item Setiap nilai $k$ menghasilkan satu kandidat plainteks.
        \item Kandidat yang bermakna menunjukkan kunci yang benar.
        \item Tidak ada kunci rahasia tambahan selain $k$.
    \end{itemize}
    \vspace{0.15cm}
    \begin{alertblock}{Sama seperti Caesar}
        \footnotesize
        Ruang kunci yang kecil membuat cipher transposisi ini tumbang oleh
        pencarian menyeluruh.
    \end{alertblock}
\end{frame}

% =====================================================================
\section{Super-Enkripsi dan Kelemahan Substitusi}
% =====================================================================

\begin{frame}{Super-Enkripsi (\textit{Product Cipher})}
    Super-enkripsi menggabungkan cipher substitusi dengan cipher transposisi.
    \begin{columns}[T,onlytextwidth]
        \begin{column}{0.48\textwidth}
            \begin{block}{Pembagian tugas}
                \footnotesize
                Substitusi mengacak identitas huruf; transposisi mengacak
                urutannya.
            \end{block}
        \end{column}
        \begin{column}{0.48\textwidth}
            \begin{exampleblock}{Urutan}
                \footnotesize
                Substitusi dahulu lalu transposisi, atau sebaliknya; keduanya
                harus disepakati.
            \end{exampleblock}
        \end{column}
    \end{columns}
    \vspace{0.15cm}
    \begin{itemize}
        \item Menggabungkan dua teknik menjadi satu cipher yang lebih sulit
              dibaca.
        \item Gagasan ini menjadi benih cipher modern yang bergantian
              melakukan substitusi dan permutasi.
    \end{itemize}
\end{frame}

\begin{frame}{Contoh Super-Enkripsi: Dua Tahap}
    \footnotesize
    Plainteks \code{hello world}.
    \vspace{0.05cm}
    \begin{block}{Tahap 1 --- Caesar $k=3$}
        \code{h}$=7\rightarrow(7+3)\bmod 26=10=$\code{K},
        \code{e}$=4\rightarrow7=$\code{H},
        \code{l}$=11\rightarrow14=$\code{O},
        \code{o}$=14\rightarrow17=$\code{R},
        sehingga \code{hello world} $\rightarrow$ \code{KHOORZRUOG}.
    \end{block}
    \begin{columns}[T,onlytextwidth]
        \begin{column}{0.48\textwidth}
            \begin{block}{Tahap 2 --- transposisi 4 kolom}
                \footnotesize
                10 huruf ditambah 2 huruf dummy \code{Z} menjadi 12 huruf.
                \[
                \begin{array}{cccc}
                    K&H&O&O\\
                    R&Z&R&U\\
                    O&G&Z&Z
                \end{array}
                \]
            \end{block}
        \end{column}
        \begin{column}{0.48\textwidth}
            \begin{exampleblock}{Hasil}
                \footnotesize
                Kolom dibaca dari kiri ke kanan: \code{KRO}, \code{HZG},
                \code{ORZ}, \code{OUZ}.\\
                Cipherteks: \code{KROHZGORZOUZ}.
            \end{exampleblock}
        \end{column}
    \end{columns}
\end{frame}

\begin{frame}{Membalik Super-Enkripsi}
    \begin{columns}[T,onlytextwidth]
        \begin{column}{0.56\textwidth}
            \begin{enumerate}
                \item Tulis cipherteks \code{KROHZGORZOUZ} ke grid 4 kolom.
                \item Baca baris demi baris: \code{KHOORZRUOGZZ}.
                \item Buang huruf dummy \code{Z} di akhir.
                \item Batalkan Caesar: \code{K}$=10\rightarrow7=$\code{h},
                      \code{H}$=7\rightarrow4=$\code{e}, dan seterusnya.
                \item Diperoleh kembali \code{hello world}.
            \end{enumerate}
            \begin{alertblock}{Urutan pembalikan}
                \footnotesize
                Transposisi dibatalkan lebih dahulu, baru substitusi, karena
                enkripsi menerapkan substitusi lebih dahulu.
            \end{alertblock}
        \end{column}
        \begin{column}{0.40\textwidth}
            \begin{block}{Arah proses}
                \footnotesize
                Enkripsi: substitusi $\rightarrow$ transposisi.\\
                Dekripsi: transposisi$^{-1}$ $\rightarrow$
                substitusi$^{-1}$.
            \end{block}
            \begin{exampleblock}{Kunci}
                \footnotesize
                Urutan yang dipakai harus sama dengan urutan pembalikannya.
            \end{exampleblock}
        \end{column}
    \end{columns}
\end{frame}

% =====================================================================
%  Satu frame penutup bagian: kelemahan substitusi monoalfabetik
% =====================================================================

\begin{frame}{Mengapa Substitusi Monoalfabetik Bocor}
    \footnotesize
    \begin{columns}[T,onlytextwidth]
        \begin{column}{0.58\textwidth}
            \begin{itemize}
                \item Substitusi tetap mempertahankan frekuensi: huruf yang
                      sering di plainteks tetap sering di cipherteks.
                \item Bahasa Inggris didominasi \code{E}, \code{T}, \code{A},
                      \code{O}, \code{I}, lalu \code{N}, \code{S}, \code{H}.
                \item Bahasa Indonesia punya urutan berbeda, sehingga tabel
                      acuan harus disesuaikan.
                \item Kriptanalis mencocokkan huruf tersering cipherteks dengan
                      huruf tersering bahasa.
                \item Bigram dan trigram mempertajam tebakan.
                \item Penyerang tidak mencari kunci, tetapi menebak pemetaan
                      huruf.
            \end{itemize}
        \end{column}
        \begin{column}{0.40\textwidth}
            \centering
            \bookfig[0.85\linewidth]{04}{frekuensi-huruf-inggris.png}
            \figsumber{Munir, \textit{IF4020 Kriptografi}, ITB.}
        \end{column}
    \end{columns}
    \vspace{0.05cm}
    \begin{alertblock}{Akibatnya}
        \footnotesize
        Ruang kunci $26!$ tidak menolong karena analisis frekuensi tidak
        menelusuri ruang kunci.
    \end{alertblock}
\end{frame}

% =====================================================================
\section{Rangkuman}
% =====================================================================

\begin{frame}{Rangkuman}
    \footnotesize
    \begin{itemize}
        \item Cipher substitusi mengganti identitas huruf; cipher transposisi
              memindahkan posisinya.
        \item Cipher abjad-tunggal memakai satu pemetaan tetap; tabelnya
              $26!$, tetapi frekuensi bahasa tetap bocor.
        \item Caesar: $c=(p+k)\bmod 26$; hanya 26 kunci sehingga tumbang oleh
              pencarian menyeluruh.
        \item Affine: $C\equiv mP+b\pmod{26}$; $m$ harus relatif prima dengan
              26 dan ruang kuncinya 300.
        \item Rail Fence menyusun plainteks zigzag pada $k$ baris lalu
              membacanya baris demi baris.
        \item Super-enkripsi merangkai substitusi dan transposisi; gagasan ini
              menjadi benih cipher modern.
    \end{itemize}
\end{frame}

\begin{frame}{Yang Harus Diingat}
    \begin{block}{Empat butir kunci}
        \begin{itemize}
            \item Cipher klasik bekerja pada huruf dan termasuk kriptografi
                  kunci-simetri.
            \item Substitusi menyembunyikan identitas huruf, transposisi
                  menyembunyikan urutan.
            \item Ruang kunci kecil membuat Caesar, Affine, dan Rail Fence
                  mudah dicoba seluruhnya.
            \item Struktur statistik bahasa tetap lolos pada substitusi
                  monoalfabetik.
        \end{itemize}
    \end{block}
    \vspace{0.25cm}
    \begin{exampleblock}{Pertemuan berikutnya}
        \textbf{Minggu 5 --- Kriptografi klasik II}
    \end{exampleblock}
\end{frame}

\end{document}
